\chapter{Perpetual Futures}\label{ch:m3:perpetual-futures}

At midnight, eight in the morning and four in the afternoon, UTC, the holders of crypto \omterm{def:m3:perpetual-futures:perp}{perpetual futures} pay
each other money. If the contract has traded above the spot price over the previous eight
hours, the longs pay the shorts; if below, the shorts pay the longs. The contract never expires, never
delivers and never needs to be rolled, and yet it stays within a fraction of a percent of spot, because
this periodic payment makes it expensive to hold the side that is pushing it away. The \omterm{def:m3:perpetual-futures:perp}{perpetual future}
is the crypto market's principal instrument of leverage: of all listed bitcoin futures traded in the first half of 2023,
perpetuals made up 75\% of the USD~27 billion average daily volume and 94\% of the USD~8 billion average open
interest, by the count Ackerer, Hugonnier and Jermann report. This chapter takes it apart:
the index it tracks and the \omterm{def:m3:perpetual-futures:index}{mark price} at which it is valued, the premium and the formula that turns it
into a funding payment, the three ways of denominating the contract, and the trade that collects the
funding.

\section{A future that never expires}

A dated future (One Quant Book 1, chapter 18) converges to spot because at expiry it becomes spot: the
long receives the underlying, or its cash value, at the final settlement price. Between now and expiry
the future can trade at a basis to spot, which the \omterm{def:m3:forward-curves-storage-and-convenience-yield:carry}{cash-and-carry trade} keeps near the cost of carry.
Take away the expiry and nothing forces convergence. The perpetual replaces delivery with a stream of
payments.

\begin{definition}[Perpetual future, funding rate, funding interval]\label{def:m3:perpetual-futures:perp}
A \emph{perpetual future}\index{perpetual future} is a futures contract with no expiry, marked to market
continuously, whose holders exchange periodic payments that push its price towards an index of the
underlying's spot price. The payment per unit of position value is the \emph{funding
rate}\index{funding rate}, positive when longs pay shorts; the period between payments is the
\emph{funding interval}\index{funding interval}, commonly eight hours.
\end{definition}

In this chapter $f$ denotes the \omterm{def:m3:perpetual-futures:perp}{funding rate} per interval, and a position of $q$ units of the
underlying held at a funding time pays $q\,M f$, where $M$ is the \omterm{def:m3:perpetual-futures:index}{mark price} defined below: a long pays
when $f > 0$ and receives when $f < 0$. If the perpetual trades rich to spot, the longs pay, which
invites arbitrageurs to sell the perpetual and buy spot, collecting the funding while hedged; their
selling pushes the perpetual back down. Funding plays the part that convergence plays for a dated future.

\begin{proposition}[The perpetual as a continuously rolled future]\label{prop:m3:perpetual-futures:roll}
Hold a long perpetual and a short spot position of the same size, and let the perpetual's price equal the
index at the start and the end. The position's total P\&L is minus the sum of the funding payments. A long
dated future bought at a basis $b$ to spot and held to expiry earns $-b$ relative to spot. Funding is
therefore the perpetual's basis paid in instalments: a perpetual whose funding averages $\bar f$ per interval
behaves like a future rolled at a basis of $\bar f$ per interval.
\end{proposition}

\begin{proof}
With equal prices at the two ends the price legs cancel and only the funding flows remain. For the dated
future, convergence moves it from $S(1+b)$ to $S_T$ while spot moves from $S$ to $S_T$, a difference of $-bS$.
\end{proof}

\begin{dated}{2026-09}{A perpetual's history and its default funding}\label{dat:m3:perpetual-futures:bitmex}
BitMEX's instrument data give 13 May 2016 as the listing date of its XBTUSD perpetual swap, an \omterm{def:m3:perpetual-futures:types}{inverse
contract} settled in bitcoin with an eight-hour \omterm{def:m3:perpetual-futures:perp}{funding interval}; on 24 September 2026 its \omterm{def:m3:perpetual-futures:perp}{funding rate} was
0.01\% per interval. Binance pays funding at 00:00, 08:00 and 16:00 UTC by default, with an interest
component of 0.01\% per interval (0.03\% a day) for most contracts.
\end{dated}

\section{Index price and mark price}

A perpetual needs two reference prices besides its own last trade.

\begin{definition}[Index price, mark price]\label{def:m3:perpetual-futures:index}
The \emph{index price}\index{index price} of a perpetual is a weighted average of the underlying's spot price
on several constituent venues, with rules that limit the influence of any one of them. The \emph{mark
price}\index{mark price} is the price at which the venue values open positions for unrealised P\&L, margin
and liquidation; it is built from the index and a smoothed basis rather than from the last trade.
\end{definition}

The index is what funding pushes the contract towards; the mark decides who is liquidated. Both are designed
against manipulation. A trader who could move one spot venue's price, or the perpetual's own last trade, for
a few seconds could otherwise trigger liquidations and profit from them
(\cref{ch:m3:margin-liquidation-and-loss-allocation}). An index that caps each constituent near the median of
all of them, and drops those that stop updating, costs more to move; a mark that uses a median of several
estimates ignores a single outlying trade.

\begin{dated}{2026-09}{Index and mark rules at a large venue}\label{dat:m3:perpetual-futures:mark}
Binance's USD-margined futures compute the index as a weighted sum of spot
prices on constituent exchanges, which it lists as including Binance, OKX, Coinbase, Kraken, Bybit and others,
as well as decentralised venues. A constituent whose latest price deviates by more than 3\% from the median of
all sources is capped at 1.03 or 0.97 times the median (1\% for some pairs such as BTCUSDT); a constituent that
cannot be reached or has not updated for five minutes gets zero weight. The \omterm{def:m3:perpetual-futures:index}{mark price} is the median of three
prices: the index times $1 + f \times (\text{time to next funding}/\text{funding period})$ with the last \omterm{def:m3:perpetual-futures:perp}{funding
rate} $f$; the index plus a 30-second moving average of the basis; and the contract's last price.
\end{dated}

\section{The premium index and the funding formula}

\begin{definition}[Premium index]\label{def:m3:perpetual-futures:premium}
The \emph{premium index}\index{premium index} measures how far the perpetual's order book is from the index:
with $I$ the index and $B^{\mathrm{imp}}$, $A^{\mathrm{imp}}$ the average prices at which a stated notional
(the impact notional) would be sold into the bids or bought from the asks,
\[
p = \frac{\max(0,\, B^{\mathrm{imp}} - I) - \max(0,\, I - A^{\mathrm{imp}})}{I}.
\]
\end{definition}

The premium is zero while the index lies between the impact bid and ask; positive when even the bids are above
it, negative when even the asks are below. Using impact prices rather than the best quotes makes it costly to
move with small orders. The venue samples it through the interval and averages the samples with weights that
grow over time, so that the premium near the funding time counts most: with $n$ samples $p_1,\dots,p_n$,
$\bar p = \sum_k k\,p_k \big/ \sum_k k$.

\begin{proposition}[The funding rule and its dead zone]\label{prop:m3:perpetual-futures:rule}
Let $i$ be the interest component per interval and $c$ the clamp. The \omterm{def:m3:perpetual-futures:perp}{funding rate}
\[
f = \bar p + \operatorname{clamp}(i - \bar p,\, -c,\, c)
\]
equals $i$ whenever $|\bar p - i| \le c$, equals $\bar p - c$ when $\bar p > i + c$ and $\bar p + c$ when
$\bar p < i - c$. With $i = 0.01\%$ and $c = 0.05\%$, funding is exactly $0.01\%$ for every average premium
between $-0.04\%$ and $+0.06\%$.
\end{proposition}

\begin{proof}
If $-c \le i - \bar p \le c$ the clamp returns $i - \bar p$ and $f = i$. If $i - \bar p < -c$ it returns $-c$;
if $i - \bar p > c$ it returns $c$.
\end{proof}

\begin{omfigure}
\begin{tikzpicture}
\begin{axis}[width=11cm,height=5cm,
  xlabel={average premium index, basis points per interval},ylabel={funding rate, bp},
  tick label style={font=\small},label style={font=\small},
  xmin=-30,xmax=30,ymin=-30,ymax=30,grid=major,grid style={gray!25},
  legend columns=2,legend style={font=\footnotesize,at={(0.5,-0.3)},anchor=north,draw=none,fill=none,/tikz/every even column/.append style={column sep=8pt}},
  table/col sep=comma]
\addplot[omThm,very thick] table[x=premium_bp,y=funding_bp]{figdata/markets-3/17-perpetual-futures/rule.csv};\addlegendentry{funding rate}
\addplot[gray,dashed,domain=-30:30,samples=2] {x};\addlegendentry{funding equal to premium}
\draw[omDef,thick,{Stealth[length=2mm]}-{Stealth[length=2mm]}] (axis cs:-4,6) -- (axis cs:6,6);
\node[font=\footnotesize,anchor=south,text=omDef] at (axis cs:1,6.5) {dead zone};
\end{axis}
\end{tikzpicture}

\omcaption{The funding rule of \cref{prop:m3:perpetual-futures:rule} with an interest component of 1 basis point
and a clamp of 5 basis points per interval: funding stays at 1 basis point for any average premium between $-4$ and
$+6$ basis points and follows the premium, 5 basis points closer to the interest component, outside that zone.
Data: the chapter's tutorial.}\label{fig:m3:perpetual-futures:rule}
\end{omfigure}

The dead zone has a consequence that surprises newcomers: in quiet markets the \omterm{def:m3:perpetual-futures:perp}{funding rate} is not the premium but
the interest component, 0.01\% every eight hours, which is $0.01\% \times 3 \times 365 = 10.95\%$ a year paid by
longs to shorts. The interest component is meant to reflect the difference between the interest rates of the quote
currency and the underlying, as in the forward points of One Quant Book 2; in practice it is a fixed parameter,
and for most of a calm year a short perpetual hedged with spot collects it.

\begin{dated}{2026-09}{The funding formula at a large venue}\label{dat:m3:perpetual-futures:formula}
Binance's \omterm{def:m3:perpetual-futures:perp}{funding rate} is $[\bar p + \operatorname{clamp}(i - \bar p, 0.05\%, -0.05\%)]/(8/N)$ for an $N$-hour
interval, with $i = 0.01\%$ per interval for most contracts. The \omterm{def:m3:perpetual-futures:premium}{premium index} uses impact bid and ask prices for an
impact margin notional of 200 USDT divided by the initial margin rate at maximum leverage, sampled every 5 seconds
and averaged with linearly increasing weights (5\,760 samples in eight hours). Funding is capped and floored at
$\pm 0.75$ times the maintenance margin ratio for major contracts such as BTCUSDT and at $\pm 2\%$ for others.
\end{dated}

\begin{omfigure}
\begin{tikzpicture}
\begin{axis}[width=11cm,height=5cm,
  xlabel={day},ylabel={basis points per interval},
  tick label style={font=\small},label style={font=\small},
  xmin=0,xmax=30,ymin=-10,ymax=16,grid=major,grid style={gray!25},
  legend columns=2,legend style={font=\footnotesize,at={(0.5,-0.3)},anchor=north,draw=none,fill=none,/tikz/every even column/.append style={column sep=8pt}},
  table/col sep=comma]
\addplot[gray,thick,mark=*,mark size=0.8pt,only marks] table[x=day,y=premium_bp]{figdata/markets-3/17-perpetual-futures/month.csv};\addlegendentry{average premium}
\addplot[omThm,very thick,const plot] table[x=day,y=funding_bp]{figdata/markets-3/17-perpetual-futures/month.csv};\addlegendentry{funding rate}
\end{axis}
\end{tikzpicture}

\omcaption{Thirty synthetic days of eight-hour intervals: calm, a rally in which the perpetual trades rich, a
sell-off in which it trades cheap, calm again. In calm periods the average premium wanders inside the dead zone and
funding is pinned at 1 basis point; in the rally longs pay up to about 10 basis points an interval; in the sell-off
shorts pay. Premium samples simulated every 5 seconds; the venue's rule applied exactly. Data: the chapter's
tutorial.}\label{fig:m3:perpetual-futures:month}
\end{omfigure}

\section{Linear, inverse and quanto contracts}

A perpetual on bitcoin can be denominated in three ways, and the choice changes what the holder is exposed to.

\begin{definition}[Linear, inverse and quanto contracts]\label{def:m3:perpetual-futures:types}
A \emph{linear contract}\index{linear contract} is quoted and settled in the quote currency (a dollar or a
\omterm{def:m3:blockchains-for-traders:stable}{stablecoin}), its P\&L linear in the price: $q\,(S - E)$ for $q$ units bought at $E$. An \emph{inverse
contract}\index{inverse contract} has a fixed notional in the quote currency but is margined and settled in the
underlying coin, so that $N$ contracts of one dollar each bought at $E$ earn $N\,(1/E - 1/S)$ coins. A
\emph{quanto contract}\index{quanto contract} is quoted on one asset but margined and settled in another at a fixed
rate per unit of price move, so that its P\&L in the settlement asset is $N\,m\,(S - E)$ for a multiplier $m$,
whatever the settlement asset's price.
\end{definition}

\begin{dated}{2026-09}{Contract terms at BitMEX}\label{dat:m3:perpetual-futures:terms}
BitMEX's instrument data on 24 September 2026 describe XBTUSD as an inverse perpetual, one US dollar of bitcoin per
contract, settled in bitcoin; ETHUSD, listed on 1 August 2018, as a quanto perpetual quoted in dollars and settled in
bitcoin with a multiplier of 100 satoshis (0.000001 bitcoin) per contract per dollar of the ether price; and XBTUSDT,
listed in November 2021, as a linear perpetual settled in USDT. All three had an eight-hour \omterm{def:m3:perpetual-futures:perp}{funding interval}.
\end{dated}

\begin{proposition}[The inverse contract's P\&L in coin]\label{prop:m3:perpetual-futures:inverse}
The P\&L in coin of a long inverse position of $N$ one-dollar contracts, $N(1/E - 1/S)$, is concave in $S$: its
derivative $N/S^2$ falls as the price rises. A long gains at most $N/E$ coins however high the price goes and loses
without bound as $S \to 0$; the short's P\&L is the mirror image. In dollars, $N(1/E - 1/S)\,S = N(S/E - 1)$ is
linear.
\end{proposition}

\begin{proof}
Differentiate twice: $N/S^2 > 0$ and $-2N/S^3 < 0$. The limits follow from $1/S \to 0$ and $1/S \to \infty$.
\end{proof}

\begin{omfigure}
\begin{tikzpicture}
\begin{axis}[width=11cm,height=5cm,
  xlabel={bitcoin price, USD},ylabel={P\&L, BTC},
  tick label style={font=\small},label style={font=\small},
  xmin=30000,xmax=120000,ymin=-1.1,ymax=1.1,grid=major,grid style={gray!25},
  xtick={30000,60000,90000,120000},xticklabels={30\,000,60\,000,90\,000,120\,000},scaled x ticks=false,
  legend columns=2,legend style={font=\footnotesize,at={(0.5,-0.3)},anchor=north,draw=none,fill=none,/tikz/every even column/.append style={column sep=8pt}},
  table/col sep=comma]
\addplot[omThm,very thick] table[x=price,y=pnl_btc]{figdata/markets-3/17-perpetual-futures/inverse.csv};\addlegendentry{long 60\,000 inverse contracts}
\addplot[omDef,thick,dashed] table[x=price,y=tangent_btc]{figdata/markets-3/17-perpetual-futures/inverse.csv};\addlegendentry{tangent at entry}
\end{axis}
\end{tikzpicture}

\omcaption{P\&L in bitcoin of a long position of 60\,000 one-dollar \omterm{def:m3:perpetual-futures:types}{inverse contracts} entered at USD~60\,000
(\cref{prop:m3:perpetual-futures:inverse}): a fall to 30\,000 costs a whole bitcoin, a rise to 120\,000 gains only
half of one. The curve lies below its tangent: a long inverse position is short convexity in coin. Data: the
chapter's tutorial.}\label{fig:m3:perpetual-futures:inverse}
\end{omfigure}

The \omterm{def:m3:perpetual-futures:types}{inverse contract} suits a holder who thinks in coin: a miner or a fund whose books are in bitcoin can hedge its
dollar value by selling \omterm{def:m3:perpetual-futures:types}{inverse contracts}, since a short inverse position plus the coins posted as margin has a
constant dollar value. A trader who thinks in dollars and goes long with coin margin is doubly long: the position
and the collateral both fall together. The \omterm{def:m3:perpetual-futures:types}{quanto contract} gives exposure to ether settled in bitcoin at a fixed
number of satoshis per dollar; its dollar P\&L is the coin P\&L times the bitcoin price, so the holder is also
exposed to bitcoin and to the correlation between the two, which the funding and the price of the contract must
compensate. \omterm{def:m3:perpetual-futures:types}{Linear contracts} margined in \omterm{def:m3:blockchains-for-traders:stable}{stablecoins} avoid both effects, and carry instead the \omterm{def:m3:blockchains-for-traders:stable}{stablecoin}'s own risk
(\cref{ch:m3:blockchains-for-traders}).

\section{Cash-and-carry and funding arbitrage}

\begin{definition}[Funding arbitrage]\label{def:m3:perpetual-futures:funding-arb}
\emph{Funding arbitrage}\index{funding arbitrage} is holding a perpetual against an offsetting position in the
underlying (spot, or another venue's perpetual) so that price moves cancel and the position earns the funding
difference; most often long spot and short perpetual to collect positive funding.
\end{definition}

It is the \omterm{def:m3:forward-curves-storage-and-convenience-yield:carry}{cash-and-carry trade} of One Quant Book 1, chapter 21, with the carry paid every eight hours instead of at
expiry. Its return depends on how much funding there is to collect. Schmeling, Schrimpf and Todorov of the BIS find
that crypto carry, the difference between futures and spot prices, averages above 10\% a year, far more than carry
in equities, bonds, currencies or commodities, is extremely volatile, reaching 60\% a year, and is highest when small
investors chase leveraged upside and arbitrage capital is scarce; high carry predicts subsequent crashes.

\begin{dated}{2026-09}{Funding on a major perpetual, by year}\label{dat:m3:perpetual-futures:years}
From Binance's public funding history for the BTCUSDT perpetual (computed by the author on 24 September 2026), the
mean \omterm{def:m3:perpetual-futures:perp}{funding rate} annualised was 17.2\% in 2020, 30.6\% in 2021, 4.2\% in 2022, 7.9\% in 2023, 11.9\% in 2024, 5.1\%
in 2025 and 2.9\% in 2026 to date; the share of negative funding events was highest in 2022 (22\%) and 2026 (26\%).
The largest single rate was 0.3\% per interval in 2020; since 2025 funding has not exceeded the default 0.01\%.
\end{dated}

\begin{omfigure}
\begin{tikzpicture}
\begin{axis}[width=11cm,height=5cm,ybar,bar width=10pt,
  xlabel={year},ylabel={\% a year},
  tick label style={font=\small},label style={font=\small},
  xmin=2019.4,xmax=2026.6,ymin=0,ymax=35,grid=major,grid style={gray!25},
  xtick={2020,2021,2022,2023,2024,2025,2026},xticklabel style={/pgf/number format/1000 sep=},
  legend columns=2,legend style={font=\footnotesize,at={(0.5,-0.3)},anchor=north,draw=none,fill=none,/tikz/every even column/.append style={column sep=8pt}},
  table/col sep=comma]
\addplot[fill=omThm!60,draw=omThm,area legend] table[x=year,y=funding_pct]{figdata/markets-3/17-perpetual-futures/years.csv};\addlegendentry{mean funding, annualised}
\addplot[fill=omDef!60,draw=omDef,area legend] table[x=year,y=trade_pct]{figdata/markets-3/17-perpetual-futures/years.csv};\addlegendentry{funding trade, return on capital}
\end{axis}
\end{tikzpicture}

\omcaption{Mean annualised funding of the BTCUSDT perpetual on Binance by calendar year (2026 to 24 September), and the
return on capital of long spot and short perpetual held a year at a constant price with 0.30\% of round-trip fees and
30.65\% of margin for the short leg (\cref{prop:m3:perpetual-futures:margin}). Data: summary statistics computed from
Binance's public funding-rate history; the trade: the chapter's tutorial.}\label{fig:m3:perpetual-futures:years}
\end{omfigure}

The trade has three weaknesses. Funding can turn negative, as it did one interval in five in 2022, and the hedged
position then pays. The short perpetual is margined separately from the spot, so a rally that the spot gains on
drains the short's margin, and if the margin is not topped up in time the short is liquidated and the position is
left unhedged at the worst moment.

\begin{proposition}[Margin to survive a rally]\label{prop:m3:perpetual-futures:margin}
A short perpetual of notional $Q$ entered at price $E$, with a maintenance margin rate $m$, survives a rally of $x$
(the price rising to $E(1+x)$) without liquidation only if its margin is at least $Q\,[x + m(1 + x)]$: for a 30\%
rally and $m = 0.5\%$, 30.65\% of the initial notional.
\end{proposition}

\begin{proof}
At $E(1+x)$ the short has lost $Qx$ and must still hold maintenance margin of $m$ times its notional, now $Q(1+x)$.
\end{proof}

The third weakness is the venue: the spot and the margin sit on exchanges, and the trade's return is a credit spread
over what a failed exchange would repay (\cref{ch:m3:centralised-exchanges}). Venues that let spot holdings count as
collateral for the short (portfolio or cross margin, \cref{ch:m3:margin-liquidation-and-loss-allocation}) remove most
of the margin drag, at the cost of concentrating both legs on one venue.

\section{Tutorial: funding, and the inverse contract}\label{tut:m3:perpetual-futures:funding}

\begin{tutorial}
\textbf{Goal.} Implement the \omterm{def:m3:perpetual-futures:premium}{premium index} and the funding rule, apply them to a simulated month of premium samples,
and plot the \omterm{def:m3:perpetual-futures:types}{inverse contract}'s P\&L in coin. \textbf{End state:}
\cref{fig:m3:perpetual-futures:rule,fig:m3:perpetual-futures:month,fig:m3:perpetual-futures:inverse} and the numbers
of the weekend problem.
\begin{tutsteps}
\item \textbf{The formula.} Premium, time-weighted average and funding, as the venue publishes them.
\omcode{code/firm/perp/firm_perp.py}{42}{56}{The premium index, its time-weighted average and the funding rate.}\label{lst:m3:perpetual-futures:funding}
\item \textbf{The month.} \texttt{month\_of\_funding(regimes())} simulates 5\,760 premium samples an interval.
\item \textbf{The trade.} The margin a short needs and the return of the funding trade.
\omcode{code/markets-3/17-perpetual-futures/python/m3_perp.py}{48}{58}{Margin to survive a rally and the funding trade's return on capital.}\label{lst:m3:perpetual-futures:trade}
\item \textbf{Run} \texttt{fig\_perp.py} for the chart data.
\end{tutsteps}
\textbf{What to change next.} Let funding be paid every hour instead of every eight and see how the dead zone scales;
charge the trade the funding of 2022 with its negative intervals; value the \omterm{def:m3:perpetual-futures:types}{quanto contract}'s correlation exposure by
simulating bitcoin and ether together.
\end{tutorial}

\section{Build: the perpetual contract engine}\label{bld:m3:perpetual-futures:perp}

\begin{build}
\bfield{Purpose} Everything the miniature firm does with perpetuals (quoting, hedging, the funding trade, the
liquidation engine of \cref{ch:m3:margin-liquidation-and-loss-allocation}) needs the venue's index, mark and funding
reproduced exactly, and P\&L in the right currency for each contract type.
\bfield{Interface} \texttt{index\_price(sources, cap, stale\_s)}; \texttt{impact\_price(levels, notional)};
\texttt{premium\_index(impact\_bid, impact\_ask, index)}; \texttt{average\_premium(samples)};
\texttt{funding\_rate(avg\_premium, interest, clamp, cap)}; \texttt{mark\_price(index, last\_funding,
frac\_to\_funding, basis\_ma, last\_trade)}; \texttt{funding\_payment}; \texttt{pnl\_linear}, \texttt{pnl\_inverse},
\texttt{pnl\_quanto}.
\bfield{Rules} Parameters per venue and contract, read from its documentation and versioned; no source, no index (an
error, never a stale value); rates per interval.
\bfield{Acceptance tests} \texttt{code/firm/perp/tests/}: capping of an outlier and exclusion of a stale source;
impact price through two levels; premium zero inside the impact spread; funding inside and outside the dead zone and at
the cap; the median mark; P\&L of the three contract types.
\bfield{Stretch} Hourly and dynamic \omterm{def:m3:perpetual-futures:perp}{funding intervals}; reproduction of a venue's published funding history from its
order-book data; the quanto correlation adjustment.
\end{build}

\begin{omsources}
\item Binance, ``Introduction to Binance Futures Funding Rates'' and ``How to Calculate the Mark Price and the Price
Index in USDS-M Futures'', support pages; public funding-rate endpoint. Accessed
September 2026.
\item BitMEX, instrument endpoint of the public API (XBTUSD, ETHUSD, XBTUSDT), 24 September 2026.
\item M.~Schmeling, A.~Schrimpf and K.~Todorov, ``Crypto carry'', BIS Working Papers 1087, April 2023.
\item D.~Ackerer, J.~Hugonnier and U.~Jermann, ``Perpetual Futures Pricing'', working paper, 27 February 2025.
\end{omsources}

\section{Exercises}

\begin{exercise}[$\star$]\label{exo:m3:perpetual-futures:1}
A long position of 2 bitcoin in a linear perpetual is held through a funding time with the mark at 60\,000 and a
\omterm{def:m3:perpetual-futures:perp}{funding rate} of 0.03\%. Who pays whom, and how much?
\end{exercise}

\begin{exercise}[$\star$]\label{exo:m3:perpetual-futures:2}
With an interest component of 0.01\% and a clamp of 0.05\%, what is the \omterm{def:m3:perpetual-futures:perp}{funding rate} for average premiums of 0.02\%,
0.10\% and $-0.08\%$?
\end{exercise}

\begin{exercise}[$\star$]\label{exo:m3:perpetual-futures:3}
Why does the venue use a median of several prices for the mark rather than the last trade?
\end{exercise}

\begin{exercise}[$\star\star$]\label{exo:m3:perpetual-futures:4}
A trader is long 60\,000 \omterm{def:m3:perpetual-futures:types}{inverse contracts} at 60\,000. What is the P\&L in bitcoin if the price goes to 50\,000? To
75\,000? In dollars?
\end{exercise}

\begin{exercise}[$\star\star$]\label{exo:m3:perpetual-futures:5}
The ETHUSD quanto pays 0.000001 bitcoin per contract per dollar of the ether price. A trader is long 10\,000 contracts
from 3\,000 to 3\,300; bitcoin is at 60\,000. What is the P\&L in bitcoin and in dollars?
\end{exercise}

\begin{exercise}[$\star\star$]\label{exo:m3:perpetual-futures:6}
Funding stays at 0.01\% per eight hours for a year. What does a short collect on USD~10 million?
\end{exercise}

\begin{exercise}[$\star\star\star$]\label{exo:m3:perpetual-futures:7}
\emph{Coding.} With \texttt{funding\_trade} and \texttt{yearly\_funding}, compute the funding trade's return on
capital each year. In which years is it below a 5\% money-market yield?
\end{exercise}

\begin{exercise}[$\star\star\star$]\label{exo:m3:perpetual-futures:8}
\emph{Find the flaw.} ``Long spot, short perpetual is market-neutral, so it cannot lose money.''
\end{exercise}

\section{Problem: The Funding Trade}

\begin{problem}[{Weekend problem --- collecting a year of funding}]\label{pb:m3:perpetual-futures:1}
A fund buys USD~100 million of bitcoin spot and sells the same notional of a linear perpetual, paying taker fees of
0.10\% on spot and 0.05\% on the perpetual at entry and at exit. It holds for a year at the funding of 2024
(\cref{dat:m3:perpetual-futures:years}); the maintenance margin rate is 0.5\%.

\medskip\noindent\textbf{Part I --- The income.}
\begin{enumerate}
\item What funding does the short collect over the year, assuming a constant price?
\item What are the fees for entering and exiting?
\item What does the default funding of 0.01\% per interval come to over a year?
\item Why was 2024's mean funding close to that default?
\item What would the income have been at 2022's funding?
\end{enumerate}

\medskip\noindent\textbf{Part II --- The margin.}
\begin{enumerate}[resume]
\item What margin must the short leg hold to survive a 30\% rally?
\item What is the capital of the trade, and its return on capital?
\item What is the return on capital if the venue counts the spot as collateral for the short and no extra margin is
needed?
\item What happens if the price rises 40\% and the margin is not topped up?
\item Why does the spot's gain in the rally not help, unless the venue cross-margins the two?
\end{enumerate}

\medskip\noindent\textbf{Part III --- The risks.}
\begin{enumerate}[resume]
\item What happens to the trade when funding turns negative?
\item What is the trade's exposure to the venue?
\item What is its exposure to the \omterm{def:m3:blockchains-for-traders:stable}{stablecoin} in which the perpetual is margined?
\item How does the basis at exit affect the result?
\item What does the BIS finding that high carry predicts crashes mean for this trade?
\end{enumerate}

\medskip\noindent\textbf{Part IV --- Judgement.}
\begin{enumerate}[resume]
\item Is the funding trade an arbitrage?
\item Why does funding exist at all, from the point of view of those who pay it?
\item How would you size the trade across several venues?
\item State the \emph{named result}: the funding trade's return on capital over 2024 with 30.65\% of margin, and the
margin needed to survive a 30\% rally.
\item In one sentence: what does a funding payment buy the long?
\end{enumerate}
\end{problem}

\section{Interview questions}

\begin{interviewq}[$\star$ \iqroles{trader}]\label{iq:m3:perpetual-futures:1}
What is a \omterm{def:m3:perpetual-futures:perp}{perpetual future} and what keeps it close to spot?
\end{interviewq}

\begin{interviewq}[$\star$ \iqroles{developer}]\label{iq:m3:perpetual-futures:2}
Why do venues liquidate on the \omterm{def:m3:perpetual-futures:index}{mark price} rather than the last trade, and how is the mark built?
\end{interviewq}

\begin{interviewq}[$\star\star$ \iqroles{researcher, trader}]\label{iq:m3:perpetual-futures:3}
Explain the difference between linear, inverse and quanto perpetuals, and who would use each.
\end{interviewq}

\begin{interviewq}[$\star\star$ \iqroles{trader}]\label{iq:m3:perpetual-futures:4}
Funding has been 0.01\% for weeks while the perpetual trades a few basis points above the index. Why?
\end{interviewq}

\begin{interviewq}[$\star\star$ \iqroles{risk}]\label{iq:m3:perpetual-futures:5}
What can go wrong with a long-spot, short-perpetual book of USD~500 million?
\end{interviewq}

\begin{interviewq}[$\star\star\star$ \iqroles{developer, researcher}]\label{iq:m3:perpetual-futures:6}
You must predict the next \omterm{def:m3:perpetual-futures:perp}{funding rate} two hours before the funding time. How?
\end{interviewq}
